% This package uses a different parsing argument for PDF and HTML output
\usepackage[pdftex,bookmarks=true,breaklinks=true,colorlinks=true,linkcolor=MidnightBlue,bookmarksopen=true,bookmarksopenlevel=1,bookmarksdepth=1]{hyperref}

\ExplSyntaxOn

\setcounter{secnumdepth}{0}

%make bigger section headers
% \titleformat*{\section}{\hspace*{-0.5em}\LARGE\bfseries}
\titleformat{\section}
    {\LARGE\bfseries}   % The style of the section title
    {Part \arabic{section}:}                             % a label prefix
    {0.75em}                          % How much space exists between the prefix and the title
    {\hspace{-0.5em}}    % A code representing the section
\titleformat{\subsection}
    {\large\bfseries}   % The style of the section title
    {\arabic{subsection}.}                             % a label prefix
    {0.75em}                          % How much space exists between the prefix and the title
    {\hspace{-0.5em}\vspace{-0.75em}}    % A code representing the section
    [\hspace{-0.5em}\titlerule]
\titleformat{\subsubsection}
    {\large\bfseries}
    {}
    {1em}
    {\hspace{-0.5em}}

% The HTML versions of mdframed environments need dumb hacks to fix parsing issues
\newmdenv[
    style=spellcontent,
    leftline=true,
    topline=true,
    rightline=true,
    bottomline=true,
]{spellcontent}


% The HTML versions of mdframed environments need dumb hacks to fix parsing issues
\newmdenv[
    backgroundcolor=LightCyan,
    style=colorenv,
    leftline=true,
    topline=true,
    rightline=true,
    bottomline=true,
    skipabove=4pt,
    innertopmargin=0.5em,
    innertopmargin=0.25em,
    innerbottommargin=0.25em,
    innerrightmargin=0.25em,
    leftmargin=0em,
]{monsterstatisticsmdframed}

\DeclareDocumentEnvironment{monsterstatistics}{}
{
  \begin{monsterstatisticsmdframed}
  \begin{fakehang}
}
{
  \end{fakehang}
  \end{monsterstatisticsmdframed}
}

\newmdenv[
    style=colorenv,
    backgroundcolor=LightCyan,
    topline=true,
]{spelltargetinginfo}

\newmdenv[
    style=colorenv,
    backgroundcolor=LightCyan,
]{augmenttargetinginfo}

\newmdenv[
    style=colorenv,
    backgroundcolor=Lavender,
]{spelleffects}

\newmdenv[
    style=colorenv,
    backgroundcolor=Lavender,
]{augmenteffects}

\newmdenv[
    style=colorenv,
    backgroundcolor=Gainsboro,
    leftline=true,
    rightline=true,
]{spellfooter}

\newmdenv[
    style=colorenv,
    backgroundcolor=Gainsboro,
    leftline=true,
    rightline=true,
    bottomline=true,
    topline=true,
]{monsterfootermdframed}
\DeclareDocumentEnvironment{monsterfooter}{}
{
    \begin{monsterfootermdframed}
    \begin{fakehang}
}
{
    \end{fakehang}
    \end{monsterfootermdframed}
}

% The HTML version of spelltwocol has to put its tabularx inside of a figure
\DeclareDocumentCommand{\spelltwocol}{m m}
{
    \par\noindent
    \begin{tabularx}{\columnwidth}{@{}~>{\lcol}X~>{\rcol}X~@{}}
        \renewcommand{\arraystretch}{0.5}
        #1 & #2
    \end{tabularx}
}

% The HTML version of threecol has to put its tabularx inside of a figure
\DeclareDocumentCommand{\sectiontwocol}{m m}
{
    \begin{tabularx}{\columnwidth}{@{}~>{\lcol}X~>{\lcol}X~@{}}
        \renewcommand{\arraystretch}{0.5}
        #1 & #2
    \end{tabularx}
}

% The HTML version of termlink can't use mbox
% args:
%   #1: reference name of term, if different from display name
%   #2: display name of term
%   #3: hyperlink prefix
\DeclareDocumentCommand{\termlink}{o m m}
{
    \IfValueTF{#1}
    {
        % Enable this to test for missing glossary definitions
        \pref{#1}
        % Enable this to test for unnecessary glossary entries
        % \label{#1}

        % TODO: remove \mbox before making major releases;
        % it fixes bizarre page break errors but makes text uglier
        \hyperlink{#3:#1}{\mbox{#2}}
    }
    {
        % Enable this to test for missing glossary definitions
        % \pref{#2}
        % Enable this to test for unnecessary glossary entries
        % \label{#2}

        \hyperlink{#3:#2}{\mbox{#2}}
    }
}

% The HTML versions of pref/pcref don't want to reference pages at all
\DeclareDocumentCommand{\pcref}{o m}
{
    \IfValueTF{#1}
    {
        #2,~page~\pageref{#1}
    }
    {
        #2,~page~\pageref{#2}
    }
}
\newcommand{\pref}[1]{page~\pageref{#1}}
\newcommand{\tref}[1]{Table~\ref{cap:#1}:~#1,~\pref{cap:#1}}
\newcommand{\trefnp}[1]{Table~\ref{cap:#1}:~#1}
\newcommand{\featpref}[1]{\pageref{feat:#1}}
\newcommand{\itempref}[1]{\pageref{item:#1}}
\newcommand{\featpcref}[1]{#1,~page~\pageref{feat:#1}}

\DeclareDocumentEnvironment{longcolumn}{}
{\onecolumn}
{\twocolumn}

\DeclareDocumentEnvironment{longtablewrapper}{}
{
    \small
    \rowcolors{1}{white}{tbrown}
    \renewcommand*{\arraystretch}{1.2}
    \setlength\LTleft{0pt}
    \setlength\LTright{0pt}
}
{}

% This allows up to a one-column page of text before starting a longtable.
% It is better than using `longcolumn` when you want to have a bit of preamble,
% but you don't want to have a full page break before the longtable starts.
\DeclareDocumentEnvironment{longtablepreface}{}
{\begin{indentedminipage}{0.49\textwidth}}
{\end{indentedminipage}}

\DeclareDocumentCommand{\sparklespaceless}{}{\includegraphics[width=0.75\fsize]{emoji/sparkles}}
\DeclareDocumentCommand{\poisonspaceless}{}{\includegraphics[width=0.8\fsize]{emoji/poison}}
\DeclareDocumentCommand{\potionspaceless}{}{\includegraphics[width=0.9\fsize]{emoji/potion}}

\ExplSyntaxOff
